\subsection{Model Design}
\label{sec:model-design}

Section~\ref{sec:bg-ml} introduced the four candidate learners at a conceptual level. We now specify their concrete configurations used in the empirical evaluation of Section~\ref{sec:experiments}. All models consume the same $78$-dimensional feature vector produced by CICFlowMeter and emit a categorical prediction over $C = 13$ classes.

\subsubsection*{Random Forest (RF).}
Scikit-learn \texttt{RandomForestClassifier} with \texttt{n\_estimators}=200, \texttt{max\_features}=$\sqrt{d}$, \texttt{class\_weight}=\texttt{balanced\_subsample}, unlimited depth. The \texttt{balanced\_subsample} option re-weights each bootstrap sample so that minority classes contribute proportionally without double-counting SMOTE-generated examples.

\subsubsection*{XGBoost.}
\texttt{XGBClassifier} with \texttt{objective}=\texttt{multi:softprob}, \texttt{n\_estimators}=500, \texttt{max\_depth}=8, \texttt{learning\_rate}=0.1, \texttt{subsample}=0.8, \texttt{colsample\_bytree}=0.8, L2 regularisation $\lambda=1.0$, histogram split finding (\texttt{tree\_method}=\texttt{hist}). Class imbalance is handled via sample weight vector $w_i = N / (C \cdot n_{y_i})$ applied in addition to SMOTE.

\subsubsection*{LightGBM.}
\texttt{LGBMClassifier} with \texttt{n\_estimators}=500, \texttt{max\_depth}=8, \texttt{learning\_rate}=0.1, \texttt{num\_leaves}=63, \texttt{is\_unbalance}=\texttt{True}. LightGBM's leaf-wise growth strategy and histogram-based binning make it significantly faster than XGBoost on wide feature matrices and more sensitive to synthetic samples near class boundaries---an advantage on SMOTE-augmented data.

\subsubsection*{Decision Tree (DT).}
Scikit-learn \texttt{DecisionTreeClassifier} with unlimited depth, Gini impurity criterion, and \texttt{class\_weight}=\texttt{balanced}. It serves as a fast, interpretable lower bound to quantify the added value of ensembling.

\subsubsection*{Why Four Models?}
The four learners allow three targeted comparisons: (i) single tree vs.\ bagging (DT vs.\ RF) to isolate ensemble gain; (ii) bagging vs.\ boosting (RF vs.\ XGBoost) to compare aggregation strategies; and (iii) level-wise vs.\ leaf-wise boosting (XGBoost vs.\ LightGBM) to assess the effect of the growth strategy under SMOTE-augmented training data.
